%=====================================================================
\chapter{Relationships and Causality}\label{ch:causality}
%=====================================================================
\locator{2}{Part II --- From Data to Insight}

\begin{outcomes}
By the end of this chapter you will be able to:
\begin{tight}
  \item \textbf{Explain} why correlation does not imply causation, and name the four alternative explanations for any observed association.
  \item \textbf{Identify} confounders, colliders and mediators in a described scenario.
  \item \textbf{Design} a randomised controlled experiment (A/B test) with a defensible sample size.
  \item \textbf{Select} an appropriate quasi-experimental design when randomisation is impossible.
  \item \textbf{Distinguish} the questions prediction can answer from those that require causal inference.
\end{tight}
\end{outcomes}

\begin{prereq}
\chref{eda} (Simpson's paradox, which is confounding seen in a chart) and
\chref{inference} (hypothesis testing, effect size, power).
\end{prereq}

\begin{keyterms}
causation $\bullet$ confounder $\bullet$ collider $\bullet$ mediator $\bullet$
reverse causation $\bullet$ spurious correlation $\bullet$ randomised controlled
trial $\bullet$ A/B test $\bullet$ control group $\bullet$ counterfactual $\bullet$
natural experiment $\bullet$ difference-in-differences $\bullet$ selection bias
\end{keyterms}

\section{The Four Explanations for Any Correlation}

You observe that students who attend the optional revision session score higher.
Before concluding the session helps, note that there are four possible
explanations, and only one of them is the one you want.

\dsfig{%
\begin{tikzpicture}[font=\scriptsize,
  v/.style={circle, draw=#1, fill=#1!10, text=#1!50!black, minimum size=1.0cm,
            font=\scriptsize\bfseries, inner sep=1pt},
  ce/.style={-{Stealth[length=2.2mm]}, line width=1.1pt},
  ttl/.style={font=\scriptsize\bfseries, align=center}]

% --- 1. causation
\begin{scope}[shift={(0,0)}]
  \node[ttl, text=greenhl!50!black] at (1.1,1.6) {1. Real causation\\{\tiny (what you hoped)}};
  \node[v=greenhl] (x1) at (0,0) {X};
  \node[v=greenhl] (y1) at (2.2,0) {Y};
  \draw[ce, greenhl] (x1) -- (y1);
\end{scope}

% --- 2. reverse
\begin{scope}[shift={(3.9,0)}]
  \node[ttl, text=goldhl!50!black] at (1.1,1.6) {2. Reverse causation\\{\tiny Y causes X}};
  \node[v=goldhl] (x2) at (0,0) {X};
  \node[v=goldhl] (y2) at (2.2,0) {Y};
  \draw[ce, goldhl] (y2) -- (x2);
\end{scope}

% --- 3. confounding
\begin{scope}[shift={(7.8,0)}]
  \node[ttl, text=accent!70!black] at (1.1,1.6) {3. Confounding\\{\tiny Z causes both}};
  \node[v=accent] (x3) at (0,0) {X};
  \node[v=accent] (y3) at (2.2,0) {Y};
  \node[v=accent] (z3) at (1.1,1.0) {Z};
  \draw[ce, accent] (z3) -- (x3);
  \draw[ce, accent] (z3) -- (y3);
  \draw[dashed, gray!60] (x3) -- (y3);
\end{scope}

% --- 4. chance
\begin{scope}[shift={(11.7,0)}]
  \node[ttl, text=purplehl!60!black] at (1.1,1.6) {4. Chance\\{\tiny you tested enough things}};
  \node[v=purplehl] (x4) at (0,0) {X};
  \node[v=purplehl] (y4) at (2.2,0) {Y};
  \draw[dashed, gray!60] (x4) -- (y4);
  \node[font=\tiny, text=purplehl!60!black] at (1.1,-0.75) {see \chref{inference}};
\end{scope}
\end{tikzpicture}}%
{The four explanations for an observed association. Data alone cannot distinguish them ---
the same correlation coefficient is consistent with all four. Only a design can.}{four-explanations}

\begin{worked}[Applying the Four to One Finding]
\emph{``Students who attend revision sessions score 12 marks higher.''}

\smallskip
\textbf{1. Causation.} The session teaches them things. Plausible.

\textbf{2. Reverse causation.} Students already doing well feel they can spare an
evening; struggling students are behind on coursework and cannot. Also plausible,
and it predicts the same correlation.

\textbf{3. Confounding.} \emph{Conscientiousness} causes both attendance at
optional sessions and higher marks. The session may contribute nothing. Highly
plausible --- and it is the default explanation for almost every ``optional
activity helps'' finding in education.

\textbf{4. Chance.} If the department examined 30 optional activities, one showing
a 12-mark gap is unremarkable.

\smallskip
\textbf{The 12 marks is not wrong; the interpretation is unsupported.} If the
department cancels the session on the strength of a null finding, or expands it on
the strength of this one, both decisions rest on nothing. What settles it is
\emph{assignment}: decide who attends by a mechanism unrelated to the student.
\end{worked}

\section{Confounders, Colliders and Mediators}

Three different roles a third variable can play. They look similar in a dataset
and demand opposite treatment, which is why they are worth separating carefully.

\dsfig{%
\begin{tikzpicture}[font=\scriptsize,
  v/.style={circle, draw=#1, fill=#1!10, text=#1!50!black, minimum size=0.95cm,
            font=\scriptsize\bfseries, inner sep=1pt},
  ce/.style={-{Stealth[length=2.2mm]}, line width=1.1pt}]

% confounder
\begin{scope}[shift={(0,0)}]
  \node[font=\scriptsize\bfseries, text=accent] at (1.3,2.0) {CONFOUNDER};
  \node[v=accent] (z) at (1.3,1.15) {Z};
  \node[v=primary] (x) at (0,0) {X};
  \node[v=primary] (y) at (2.6,0) {Y};
  \draw[ce, accent] (z) -- (x);  \draw[ce, accent] (z) -- (y);
  \node[align=center, text width=4.0cm, font=\tiny, text=gray!55!black] at (1.3,-1.35)
       {Z causes both. \textbf{Control for it} --- otherwise the X--Y association is fake.
        \emph{Programme difficulty, in \dsref{simpsons}.}};
\end{scope}

% mediator
\begin{scope}[shift={(5.0,0)}]
  \node[font=\scriptsize\bfseries, text=greenhl!50!black] at (1.3,2.0) {MEDIATOR};
  \node[v=greenhl] (m) at (1.3,1.15) {M};
  \node[v=primary] (x2) at (0,0) {X};
  \node[v=primary] (y2) at (2.6,0) {Y};
  \draw[ce, greenhl] (x2) -- (m);  \draw[ce, greenhl] (m) -- (y2);
  \node[align=center, text width=4.0cm, font=\tiny, text=gray!55!black] at (1.3,-1.35)
       {X causes Y \emph{through} M. \textbf{Do not control for it} --- you would erase
        the very effect you are measuring.};
\end{scope}

% collider
\begin{scope}[shift={(10.0,0)}]
  \node[font=\scriptsize\bfseries, text=purplehl] at (1.3,2.0) {COLLIDER};
  \node[v=purplehl] (c) at (1.3,1.15) {C};
  \node[v=primary] (x3) at (0,0) {X};
  \node[v=primary] (y3) at (2.6,0) {Y};
  \draw[ce, purplehl] (x3) -- (c);  \draw[ce, purplehl] (y3) -- (c);
  \node[align=center, text width=4.0cm, font=\tiny, text=gray!55!black] at (1.3,-1.35)
       {Both cause C. \textbf{Controlling for it creates} a spurious X--Y association
        that does not exist.};
\end{scope}
\end{tikzpicture}}%
{Three roles for a third variable. ``Control for everything you have'' is not a safe
default: adjusting for a mediator hides a real effect, and adjusting for a collider
manufactures a false one.}{confounder-collider}

\begin{conceptbox}[A Collider You Will Meet]
Suppose both \emph{technical skill} and \emph{interview preparation} independently
increase the chance of being hired. Among \emph{hired} employees, the two become
negatively correlated --- someone who got in with little preparation must have been
very skilled. Studying only hired staff, you would ``discover'' that preparation
harms skill. Selecting on the outcome (hiring) is conditioning on a collider, and
it is the mechanism behind a great deal of survivorship bias.
\end{conceptbox}

\section{Experiments: The Gold Standard}

\begin{definitionbox}[Randomised Controlled Trial]
Units are assigned to treatment or control \emph{at random}. Randomisation makes
the two groups equivalent in expectation on \emph{every} variable --- measured,
unmeasured and unimagined --- so any difference in outcome can be attributed to the
treatment.
\end{definitionbox}

\begin{conceptbox}[Why Randomisation Is So Powerful]
Statistical control can only adjust for confounders you thought of and measured.
Randomisation handles the ones you never knew existed. That is the whole argument,
and it is why one small randomised study can outweigh a large observational one.
\end{conceptbox}

\dsfig{%
\begin{tikzpicture}[font=\scriptsize,
  b/.style={rounded corners=3pt, draw=#1, fill=#1!9, text=#1!50!black,
            text width=2.6cm, align=center, minimum height=0.95cm, font=\scriptsize\bfseries},
  arr/.style={-{Stealth[length=2.2mm]}, gray!60, line width=1pt}]

  \node[b=primary]   (pop) at (0,0)      {Eligible\\population};
  \node[b=secondary] (smp) at (3.3,0)    {Random\\sample};
  \node[b=greenhl]   (tr)  at (6.9,0.95) {TREATMENT\\new lab format};
  \node[b=goldhl]    (ct)  at (6.9,-0.95){CONTROL\\old format};
  \node[b=greenhl]   (o1)  at (10.3,0.95) {outcome\\$\bar{y}_T$};
  \node[b=goldhl]    (o2)  at (10.3,-0.95){outcome\\$\bar{y}_C$};
  \node[b=purplehl]  (eff) at (13.4,0)    {effect\\$= \bar{y}_T - \bar{y}_C$};

  \draw[arr] (pop) -- (smp);
  \draw[arr] (smp) -- (tr);
  \draw[arr] (smp) -- (ct);
  \draw[arr] (tr) -- (o1);
  \draw[arr] (ct) -- (o2);
  \draw[arr] (o1) -- (eff);
  \draw[arr] (o2) -- (eff);

  \node[font=\tiny\itshape, text=accent!80!black, align=center, text width=3.0cm] at (5.1,-2.15)
       {\textbf{Randomise here.} This single step is what makes the final subtraction
        a causal estimate rather than a correlation.};
  \draw[-{Stealth[length=1.8mm]}, accent!70, dashed] (5.1,-1.55) -- (5.1,-0.35);
\end{tikzpicture}}%
{The structure of a randomised experiment. Everything hinges on the random assignment step:
without it, the final difference confounds the treatment with whatever made people end up
in each group.}{rct-structure}

\subsection{Sizing an Experiment}

\begin{worked}[How Many Students Do You Need?]
You want to detect a 10 percentage-point rise in pass rate (from 68\% to 78\%) with
80\% power at $\alpha = 0.05$.

\smallskip
For two proportions, a workable approximation is
\[
n \;\approx\; \frac{16 \,\bar{p}(1-\bar{p})}{\delta^2}
\quad\text{per group, where } \delta \text{ is the difference to detect.}
\]

With $\bar{p} = 0.73$ and $\delta = 0.10$:
\[
n \approx \frac{16 \times 0.73 \times 0.27}{0.01}
        = \frac{16 \times 0.1971}{0.01}
        = \frac{3.154}{0.01} \approx 315 \text{ per group}
\]

\textbf{So 630 students in total} --- and the worked example in \chref{inference},
which had 100 per group, was therefore never capable of detecting this effect. That
is why its $p = 0.078$ told you almost nothing.

\smallskip
\textbf{Halve the effect you want to detect} ($\delta = 0.05$) and $n$ rises to
about 1{,}260 per group: the sample scales with $1/\delta^2$. Small effects are
extremely expensive to establish, which is worth knowing \emph{before} promising a
stakeholder an answer.
\end{worked}

\section{When You Cannot Randomise}

Often randomisation is impossible --- you cannot randomly assign students to be
disadvantaged, or randomly expose half your servers to attack. Quasi-experimental
designs recover part of the causal claim.

\rowcolors{2}{white}{lightbg}
\begin{longtable}{@{}L{3.5cm}L{5.2cm}L{5.8cm}@{}}
\toprule
\rowcolor{primary!15}
\textbf{Design} & \textbf{Idea} & \textbf{Key assumption (state it explicitly)} \\
\midrule
\endhead
Before/after & Compare the same units before and after a change & Nothing else changed at the same time --- usually false \\
Difference-in-differences & Compare before/after change in a treated group against an untreated one & The two groups' trends would have moved in parallel without the treatment \\
Matching & Pair each treated unit with a similar untreated one & All relevant confounders were measured and used in matching \\
Regression discontinuity & Compare units just either side of a cutoff (e.g.\ a 40\% pass mark) & Units cannot precisely manipulate which side they land on \\
Instrumental variable & Use something that shifts treatment but not the outcome directly & The instrument affects the outcome \emph{only} through the treatment \\
\bottomrule
\end{longtable}

\begin{alertbox}[Before/After Alone Is Not Evidence]
``We deployed the new firewall rules and incidents fell 40\%'' is compatible with:
the rules worked; attack volume fell for unrelated reasons; the reporting process
changed; a seasonal trough. Add a comparison group that did \emph{not} receive the
change and you have difference-in-differences --- a genuinely different quality of
claim, for very little extra work.
\end{alertbox}

\section{Prediction Does Not Need Causation --- Until It Does}

\begin{conceptbox}[The Distinction That Decides Your Method]
For \textbf{prediction}, correlation is enough. A model that flags at-risk students
using forum inactivity works whether or not inactivity \emph{causes} failure.

For \textbf{intervention}, correlation is useless. Forcing students to post on the
forum will not help if inactivity was merely a symptom. The moment anyone asks
``so what should we do?'', you have left prediction and entered causal inference
--- and Parts III and IV, which are entirely about prediction, will not answer it.
\end{conceptbox}

%---------------------------------------------------------------------
\begin{fourlenses}{Causal Claims}
\lensLA{\emph{Confounder:} prior attainment causes both engagement and final marks,
so any ``engagement helps'' finding must adjust for it. \emph{Design:} randomise
\emph{which students are offered} the intervention, not who takes it --- then
compare by offer (intention to treat), which stays causal even though uptake is
voluntary. \emph{Ethical constraint:} you may not withhold a support service
believed to work, so stepped-wedge designs, where everyone eventually receives it,
are the norm.}

\lensPM{\emph{Confounder:} project complexity causes both the choice of methodology
and the risk of overrun, so ``agile projects succeed more'' is almost always
confounded --- simple projects get agile. \emph{Design:} randomising methodology
across projects is rarely feasible; difference-in-differences across teams that
adopted a practice at different times is the practical alternative.
\emph{Collider warning:} studying only \emph{completed} projects conditions on
survival and will mislead you.}

\lensIOT{\emph{Confounder:} ambient temperature drives both cooling-fan speed and
failure rate. \emph{Advantage:} devices are far easier to randomise than people ---
push a firmware change to a random half of the fleet and you have a true RCT, with
no consent problem and thousands of units. This makes IoT one of the few domains
where clean causal evidence is cheap, and it is badly under-exploited.}

\lensSEC{\emph{Reverse causation:} do more security tools cause fewer incidents, or
do organisations that suffer incidents buy more tools? Both, which makes
cross-sectional comparisons nearly meaningless. \emph{Design:} A/B test detection
rules on randomly split traffic. \emph{Collider warning:} analysing only
\emph{detected} attacks conditions on detection, so you will systematically
conclude that attacks look like the ones your current tools already catch.}
\end{fourlenses}

%---------------------------------------------------------------------
\begin{lab}[Simulating a Confounder]
The clearest way to understand confounding is to build one and watch it fool you.

\begin{lstlisting}[language=Python]
import numpy as np, pandas as pd
rng = np.random.default_rng(7)
n = 2000

# Z = conscientiousness. It causes BOTH attendance and marks.
Z = rng.normal(0, 1, n)

attends = (Z + rng.normal(0, 0.6, n)) > 0.3        # caused by Z
mark    = 55 + 9*Z + rng.normal(0, 5, n)           # caused by Z ONLY
#          note: `attends` does NOT appear in `mark` -- zero true effect

df = pd.DataFrame({"Z": Z, "attends": attends, "mark": mark})

naive = df[df.attends].mark.mean() - df[~df.attends].mark.mean()
print(f"naive 'effect' of attending: {naive:+.2f} marks")   # ~ +10, and entirely fake

# Adjust for the confounder by comparing within strata of Z:
df["band"] = pd.qcut(df.Z, 5)
adj = (df.groupby("band", observed=True)
         .apply(lambda g: g[g.attends].mark.mean() - g[~g.attends].mark.mean(),
                include_groups=False)
         .mean())
print(f"after adjusting for Z:       {adj:+.2f} marks")     # ~ 0, the truth
\end{lstlisting}

\textbf{The lesson:} the naive comparison reports a large, highly significant
effect of a variable that has no effect at all. No amount of extra data would have
corrected it --- only the adjustment did.
\end{lab}

%---------------------------------------------------------------------
\begin{checkpoint}
\begin{tightnum}
  \item Ice cream sales correlate with drowning deaths. Name the confounder and
        classify the association using \dsref{four-explanations}.
  \item You want to know whether a support programme helps. Why is comparing
        participants with non-participants insufficient, even with a large sample?
  \item A study of hired employees finds that interview preparation correlates
        negatively with technical skill. What is the likely explanation?
\end{tightnum}
\end{checkpoint}

%---------------------------------------------------------------------
\begin{chaptersummary}
\begin{tight}
  \item Any correlation admits four explanations: causation, reverse causation,
        confounding and chance. Data alone cannot separate them.
  \item Control for confounders; do \emph{not} control for mediators or colliders
        --- adjusting for a collider creates associations that are not there.
  \item Randomisation balances confounders you never thought of, which is why an
        RCT beats a much larger observational study.
  \item Required sample size scales as $1/\delta^2$: detecting small effects is
        disproportionately expensive.
  \item When randomisation is impossible, quasi-experimental designs work --- but
        each rests on an assumption you must state.
  \item Prediction tolerates correlation; intervention requires causation. Know
        which question you are being asked.
\end{tight}
\end{chaptersummary}

\begin{reviewq}
  \item \textbf{(MCQ)} Adjusting for which type of variable \emph{creates} a
        spurious association? (a)~Confounder (b)~Mediator (c)~Collider
        (d)~Instrument.
  \item \textbf{(MCQ)} To detect an effect half as large, the required sample size
        multiplies by roughly: (a)~2 (b)~4 (c)~8 (d)~$\sqrt{2}$.
  \item \textbf{(Short)} Explain in three sentences why randomisation defeats
        confounders that were never measured.
  \item \textbf{(Short)} State the key assumption of difference-in-differences and
        give one situation where it plainly fails.
  \item \textbf{(Short)} Why is analysing only \emph{detected} attacks a collider
        problem, and what does it do to your conclusions?
  \item \textbf{(Applied)} A university reports that students using the new mobile
        app achieve 6 marks more on average. Design a study --- randomised if
        possible, quasi-experimental if not --- that would support a causal claim.
        State your sample size reasoning and the assumption your design relies on.
  \item \textbf{(Applied)} An operations team says ``since we introduced automated
        patching, incidents are down 35\%''. List three alternative explanations and
        describe the smallest change to their analysis that would make the claim
        credible.
\end{reviewq}
